\chapter{恋词茶会}

玲行走在这座喜鹊搭成的桥上，感受着脚底传来鹊儿扇动翅膀带来的微微振动。缓缓流动的银河在桥下淌过，密集而耀眼的星光遮蔽了喜鹊羽毛的微光；这可是个难见的景象。“看来相比‘沉舟侧畔千帆过’，还是‘沉舟侧畔千星过’更有味道。”玲开始她的文学性联想，“那么下一句就该是‘病树前头万宙春’了。”不过，这样的探讨人生哲理的诗句放在这里是不合适的；因为鹊桥首先是个应该让人感到幸福的象征爱情的意象。想到这里，玲忍不住哼起歌来：
\begin{quote}
皎皎明月，携我游空\\
步星为舞，沐彼春风\\
荧惑之上，岁星之东\\
执子素手，贻我丹红\par
\vspace{0.5\baselineskip}%
清歌既作，长萦心胸\\
与子同曲，永以为诵\\
彼美孟姜，德音不亡\\
缟衣如荼，永矢弗谖
\end{quote}

如果这算作\textit{Fly Me to the Moon}的翻译作品，那显然是不合格的。并且，玲现在已经行走在银河之上，再谈什么月球、木星、火星就有点太小家子气了。“应该改为‘Fly Me to the Galaxy’；中文标题就叫‘共览星河’。”玲思考着将其改造的可能，“不如‘共揽星河’；还是直接‘逐鹿星河’吧！”由玲的想法可知，玲目前并未处于恋爱关系中。

这样边走边想，玲发现不久后就从织女星走到了牛郎星（这是玲通过恒星光谱确定的）。这样的事实立刻使玲产生了一个疑问：自己刚才的速度有多快呢？玲于是开始计算。先从简单的开始；已知牛郎星距离地球\(16.7\)光年，织女星距离地球\(25.0\)光年；它们属于“夏季大三角”（有种三角恋的气息），在夜空中相隔约\(34\)度；于是，代入余弦定理：
\[
c = \sqrt{a^2 + b^2 - 2ab\cos\theta}
\]
可得两者之间的距离为\(14.5\)光年。

玲刚才走了大概十几分钟；这是她经历的固有时间；为了计算简便，不妨设她刚好走了\(1000\)秒（\(\Delta\tau = \qty{1.00e3}{s}\)）。她在这\(1000\)秒里走过了\(14.5\)光年；或者说（为了单位统一），\(d=\qty{4.57e8}{光秒}\)。根据狭义相对论的时间膨胀公式，在地球参照系中：
\[
\Delta t = \gamma \Delta\tau
\]
其中洛伦兹因子\(\gamma\)是：
\[
\gamma = \frac{1}{\sqrt{1 - (v/\mathrm{c})^2}}
\]
同时别忘了最基本的速度、时间和距离的关系：
\[
d = v \Delta t
\]
联立这三个等式，我们可以得到：
\[
\frac{v}{\mathrm{c}} = \frac{d}{\sqrt{\Delta\tau^2 \mathrm{c}^2 + d^2}}
\]
代入刚才的数值计算，最后我们可以得到：
\[
v = 0.9999999999976\mathrm{c}
\]
“这已经十分接近光速了。”玲评价道，“这样的计算结果是无效的，因为超出了原始数据的精度。”

也就是说，以地球上的观测者的视角，玲以接近光速的速度，花费了约14.5年，从织女星移动到了牛郎星。“真是天上一刻，地下十年。”不过，这样的观察与玲自己的体验是毫无关系的；玲因此想到自己应该使用的事实上是尺缩公式：
\[
L = \frac{d}{\gamma}
\]
\(\gamma\)在先前的运算中已经可以算出来了，在这里可以直接代入；最后得到：
\[
L \approx \qty{1e3}{ls} = \qty{3e11}{m}
\]
“我早该想到的，”玲评价道，“对我来说，实际上是牛郎星与我的距离急剧地缩短了。”

玲接着开始思考如何弥补自己这（在地球人看来）损失的14.5年。“这就是狭义相对论不再适用的地方了，”玲想到，“因为我的行走涉及加速和减速，以及运动方向的改变。”玲思考着如果自己走回去会怎么样，“这就是经典的双生子佯谬了。”

可惜走回去已经不再可能了，因为喜鹊们突然惊慌地飞散，扑腾的翅膀勾勒出缭乱的荧光；这座鹊桥就这样迅速地瓦解了，并不留一丝痕迹。玲赶紧扑到对岸的地面上。她这才发现了导致鹊桥解体的始作俑者：一只粉白色毛发的猫正在扑食组成鹊桥的鹊儿，它的爪子里正握着被胡乱扯下脑袋的一只可怜的受害者。

“你在干什么！”玲感到不适，“不许你破坏生态！”玲旋即发觉自己的斥责是可笑的。正常的猫可不会有保护生态的概念；即使有，这样的由过量的喜鹊构成的桥算是正常的应该保护的生态吗？（“论银河上下游的普通喜鹊种群动态及分布模式”。）另外，斥责的前提是对方听得懂自己说的话。

“它们完成了使命，”粉猫擦了擦嘴上留下的血，“而我不过是从中渔利。”它瘆人地笑了起来，嘴角划破脸颊，硕大的紫色瞳孔注视着玲。

如果一只猫会说话并且会这样笑，那么它就是刘易斯的柴郡猫了。“柴郡猫！”玲上前打招呼。

“如果‘柴郡猫’是你希望的谓词，那你恐怕要失望了，”柴郡猫依然笑着，也不知道它有没有感到冒犯，“我现在有了名字。我叫‘彩名’。”

玲一点也不喜欢这个名字。“彩”立刻让她联想到“杂多”，而“名”则使她联想到“规定”；于是“彩名”连起来便是“杂多的规定性”。“这样下去我就要被海德格尔遮蔽了。”玲感到不爽，决定继续称呼它为柴郡猫。

“你也会像柴郡猫一样身体一部分一部分地消失和出现吗？”玲对这个问题感到好奇。

“按照前定和谐，”玲认为它的意思是“按照原著的情节”，“你应当先向我问路。”

“麻烦你告诉我，”玲决定放空大脑配合它，“我应该走哪条路？”

“这取决于你要去什么地方。”猫说。

“去什么地方，我并不很在乎。”这是实话。玲心想，“或者说，我更想早点能够回家睡觉。”

“那么，你走哪条路也就无所谓了。”猫回答。

“这并不一样，”玲质疑，“如果我不能在选择前知道我有哪些选择，以及我的选择的后果，那么这样的选择就不是自为的。”

“你一定要自欺地为你的选择找一个根基吗？这样的话，你可真是个彻头彻尾的疯子。”玲无法看出柴郡猫此时的笑容是嘲讽还是怜悯。

怎么看都是你们和这个世界疯了吧，玲在心中反驳。

柴郡猫似乎听到了玲的心声。它回应道：“你经历着一系列超出你的认知能力和思维方式的事件；你把它们归为奇事怪事，以此和它们划清界限，把自己当作一个局外人。”它的眼睛仍然钉在玲身上，令玲感到不安。

“那么接下来你是不是要因为我没有在母亲的葬礼上哭泣而宣判我的死刑了。”玲想象着这个情节。不过，玲的母亲并没有被设定为已去世。

柴郡猫没有理会玲的小聪明，继续说着：“你一直在尝试用你所谓的理性及其逻辑来解释这个世界的种种你认为的荒谬之处，却从未反思你所仰赖的、认为无需任何中介的逻辑，它本身就有可能是荒谬的。”

“如果我连我思维中的逻辑都不能依赖，我还能依赖什么呢？”玲被它绕晕了。

柴郡猫噗嗤一笑。它看来是不想再继续这个话题了：“我还是给你介绍一下这附近都住着谁吧。”

“左边的路通往薛定谔猫居住的盒子。”柴郡猫举起爪子指了指，“它是我的同类……或者说，朋友。我没见过它；我猜，它既然叫‘猫’的话，那说不定和我一样是一只猫。我和它一直有通过超距通讯相互联络。”

是活着的猫还是死去的猫就不好说了。玲微微点了点头，表示自己在听。

“右边是织宇蚕的巢穴；不过，你要去的话，可别在卡拉比–丘迷宫里迷了路。”柴郡猫的两只手比划着，摆出奇怪的姿势，“它最近正愁没人承办今年的世界字符串大会。”

玲想着世界字符串大会和卡拉比–丘流形之间的联系。最后她恍然大悟：“它说的应该是世界弦论大会；这两者的英文是一样的。”

“前面是拜词教的聚会。”柴郡猫用尾巴指了指它的身后，也就是玲的前方，“我没去过、也不想去那地方；那里都是些狂热的教徒。”

“那么，我已经告诉你你拥有的选项了；你去哪儿就由你‘自为’决定喽！”柴郡猫咯咯地笑了一声，说完就消失不见了。

玲思索着。她首先排除去左边。其一是她怕她的到访会使薛定谔猫坍缩到确定的状态；如果恰好坍缩到了死亡的状态，那么她的罪责就跳进黄河也洗不清了。（不过，正义的检察官应当惩处把猫关进盒子里的始作俑者而不是到访的玲。）其二，量子力学及其后的量子场论、量子色动力学对于演算能力的要求过高了；即使玲会算，本书的作者也无法检查是否正确。

“至于织宇蚕，祝它能够顺利开办今年的弦论大会。”玲接着考虑第二个选项，“我就不去了；我怕它一展开，我的内脏会掉出来。”

于是只剩下唯一的选项。玲只好迈步前进。还没等她走出半步，柴郡猫突然又出现了。“如果你真要去那边，”看来它还有事情叮嘱玲，“我先给你几本书看看，免得你被他们误导。”

柴郡猫说完，再次消失了。这次它消失得非常缓慢，最先是咧嘴的笑容，身体一块一块地隐去，最后是尾巴梢头，只留下地上的几本书。“这正好和刘易斯的柴郡猫是相反的，”玲回忆着，“这样的话，它刚要开始消失，就立刻不能说话了。”玲想着一具不能说话的身体意味着什么，“或许我们可以以此写一本题为《言语消失之日》的小说。”

玲拾起地上的书，立刻就看完了。
